\begin{abstract}

针对“板凳龙”盘入、调头和盘出过程中的位置、速度、碰撞规避及路径优化问题，本文建立基于阿基米德螺线、矩形碰撞检测和两段相切圆弧的参数化运动模型。基于相邻把手中心间固定直线距离的几何约束，通过数值求解非线性方程递推得到223节板凳共224个把手中心在等距螺线（螺距$p=0.55$m）上的位置，并利用速度投影关系推导各把手速度。在$t=300$s时，龙头前把手位于$(4.420,2.320)$m，速度$1.0$m/s，龙尾速度为$0.9965$m/s。建立基于矩形轮廓与螺线间距的碰撞判据，采用二分法搜索得到首次碰撞终止时刻$412.474$s，此时龙头位于$(1.210,1.943)$m，全队最大速度仍为$1.0$m/s。将最小螺距问题转化为约束优化问题，求得能使龙头恰好进入直径$9$m圆形调头区域的最小螺距为$0.4503$m。对两段相切圆弧构成的S形路径，以路径总长度为目标函数，采用黄金分割法搜索最优圆弧半径比例因子，得到$k=2.219$，对应路径长度$13.621$m。最后在问题4路径上计算各把手速度放大系数，得到$\gamma_{\max}=1.593$，从而确定龙头最大行进速度为$1.256$m/s，可保证所有把手速度不超过$2$m/s。

\vspace{0.5em}
\noindent\textbf{关键词}：板凳龙；阿基米德螺线；多刚体运动学；碰撞检测；二分法；约束优化；速度投影法
\end{abstract}
